\documentclass[ignorenonframetext]{beamer}

\usepackage[utf8x,utf8]{inputenc}
\usepackage[T1]{fontenc}
\usepackage[frenchb]{babel}
\usepackage{svg}

\uselanguage{French}
\languagepath{French}
%\usepackage[latin1]{inputenc}

\usepackage{beamerseminar}
%\usepackage[accumulated]{beamerseminar}
                                % remove ``accumulated'' option
                                % for original behaviour
%\usepackage{beamerthemeclassic}
\usepackage{beamerthemeshadow}
\title{Infinite Computations in Algorithmic Randomness and Reverse Mathematics}
\institute{\hspace{1cm}\includegraphics[scale=0.25]{lacl_logo.png}\hspace{1.8cm}\includegraphics[scale=0.35]{UPEC-logo.eps}}
\date{November 22nd, 2019}
\author[Paul-Elliot Anglès d'Auriac]{Paul-Elliot Anglès d'Auriac\\ \vspace{0.75cm}Under the supervision of \\\vspace{0.2cm}Benoît Monin\\ Pierre Valarcher}

%\newcommand{\sref}[1]{SLIDE \ref{#1}}

% CHANGED: different definition of \heading
%\newcommand{\heading}[1]{\begin{center}\large\bf #1\end{center}}
\let\heading=\frametitle

% CHANGED: Commented:
%\newpagestyle{MH}%
%  {University of Guaduas, March 13, 1998\hfil\thepage}{}
%\pagestyle{MH}

\usepackage{proof}

\usepackage{amsmath}
\usepackage{stmaryrd}
\usepackage{amssymb}
\usepackage{amsthm}
\usepackage{enumerate}
\usepackage{verbatim}
\usepackage{graphicx}
\usepackage{cancel}

\usepackage{mathrsfs}
\usepackage{textpos}
\usepackage{mathtools}
%\usepackage{extpfeil}
\usepackage{color}
%\usepackage{ulsy}

% \usepackage{fullpage} 

\usepackage{bm}
\newcommand\bfSigma{\boldsymbol{\Sigma}}
\newcommand\bfPi{\boldsymbol{\Pi}}
\newcommand\bfDelta{\boldsymbol{\Delta}}


\newcommand\Lc{\mathscr{L}}
\newcommand\Mc{\mathcal{M}}
\newcommand\Nc{\mathcal{N}}
\newcommand\Pc{\mathcal{P}}
\newcommand\Ac{\mathcal{A}}
\newcommand\Bc{\mathcal{B}}
\newcommand\Cc{\mathcal{C}}
\newcommand\Uc{\mathcal{U}}
\newcommand\Fc{\mathcal{F}}
\newcommand\Ls{\mathscr{L}}
\newcommand\Ms{\mathscr{M}}
\newcommand\Ns{\mathscr{N}}
\newcommand\Ps{\mathscr{P}}
\newcommand\As{\mathscr{A}}
\newcommand\Bs{\mathscr{B}}
\newcommand\Cs{\mathscr{C}}
\newcommand\Ds{\mathscr{D}}
\newcommand\Us{\mathscr{U}}
\newcommand\Fs{\mathscr{F}}
\newcommand\Xs{\mathscr{X}}
\newcommand\Os{\mathscr{O}}
\newcommand\Oc{\mathcal{O}}
\newcommand\Ys{\mathscr{Y}}
\newcommand\Zs{\mathscr{Z}}
\newcommand\pe{\preccurlyeq}
\newcommand\ep{\succcurlyeq}
\newcommand\tp{\mathrm{tp}}
\newcommand\di{\mathrm{div}}
\newcommand\cof{\mathrm{cof}}
\newcommand\Aut{\mathrm{Aut}}
\newcommand\card{\mathrm{card}}
\newcommand\Card{\mathrm{Card}}
\newcommand\Club{\mathrm{Club}}
\newcommand\type{\mathrm{type}}
\newcommand\Ord{\mathrm{Ord}}
\newcommand\V{\mathrm{V}}
\newcommand\trcl{\mathrm{trcl}}
\newcommand\Rg{\mathrm{Rg}}

\newcommand*{\longhookrightarrow}{\ensuremath{\lhook\joinrel\relbar\joinrel\rightarrow}}

\newtheorem{proposition}[theorem]{Proposition}%[section]
\newtheorem{Question}[theorem]{Question}%[section]
\newtheorem{Paradigme}[theorem]{Paradigme}%[section]

\begin{document}

% CHANGED: Added \frame
\begin{frame}
  \maketitle          % This won't show up when \onlynotestoo is in effect.
\end{frame}
% \frame{\heading{Table des matières}
% \tableofcontents                  % This won't show up when \onlynotestoo is in effect.
%}

\begin{frame}
  \frametitle{Quelque clés pour suivre}
  \begin{definition}
    \begin{itemize}
    \item Un calcul, c'est ...
    \item Un ensemble est calculable si ...
    \item Une fonction est calculable, si ...
    \end{itemize}
  \end{definition}
\end{frame}

\begin{frame}
  \frametitle{De forts liens entre définissabilité et calculabilité}
  \begin{itemize}
  \item Un algorithme définit un objet,
  \item Réciproquement, une définition induit un algorithme
  \end{itemize}
  Ces liens permettent d'utiliser le paradigme ``calcul'' dans d'autres domaine de la logique :
  \begin{itemize}
  \item La théorie de la preuve, qui étudie du point de vue mathématiques la prouvabilité,
  \item La théorie descriptive des ensembles, qui étudie les propriétés des ensembles de rééls définissables
  \end{itemize}
\end{frame}
\begin{frame}
  Dans cette thèse, nous nous intéressons à un type particulier de calcul : la calcul à temps infini.
\end{frame}
\begin{frame}
  Randomness:
  The study of the objects taken at random, and their properties.

  Taken at random? Informally:
  \begin{enumerate}
  \item each digit taken at random, one after another, or
  \item taken uniformly inbetween 0 and 1
  \end{enumerate}
  Some properties may happen, but have probability 0 to happen :
  \begin{enumerate}
  \item having only 0s
  \item having more 1s than 0s
  \end{enumerate}
\end{frame}
\begin{frame}
  We can use Infinite Time Turing Machines and the constructibles to define notions of randomness:
  \begin{definition}
    \begin{enumerate}
    \item ITTM-randomness,
    \item randomness over $L_\alpha$,
    \item ITTM-ML-randomness...
    \end{enumerate}
  \end{definition}
  \begin{theorem}
    Le theoreme de Hindman
  \end{theorem}
\end{frame}
\begin{frame}
  Résultats et questions
\end{frame}
\begin{fame}
Hindman's Theorem  
\end{fame}
\end{document}

